%! Author = sriadityavedantam
%! Date = 3/25/26

\section{Uniform Convergence of a Sequence of Functions}

\lesson{38}{Monday 10 November 2025 9:10}{Day 38}

\begin{definition}[Pointwise Convergence]
    Given functions $f_n, f_0 : A\to\rr$, we say $(f_n)$ converges pointwise to $f_0$ if for every $x\in A$,
    \[
        f_n(x)\to f_0(x).
    \]
\end{definition}

\begin{problem}{
    \[
        f_n(x) = x^n, \quad A = [0,1].
    \]
}
    \[
        x^n \to f_0(x) \coloneqq
        \begin{cases}
            0, & 0 \leq x < 1,\\
            1, & x = 1.
        \end{cases}
    \]
    So convergence is pointwise on $[0,1]$.
\end{problem}

\begin{problem}{
    \[
        f_n(x) = x^{1/n}, \quad A = [0,\infty).
    \]
}
    \[
        x^{1/n} \to g_0(x) \coloneqq
        \begin{cases}
            0, & x = 0,\\
            1, & x > 0.
        \end{cases}
    \]
    So convergence is pointwise on $[0,\infty)$.
\end{problem}

\begin{problem}{
    \[
        f_n(x) = \dfrac{x^2 + nx}{n}, \quad A = \rr.
    \]
}
    \[
        \dfrac{x^2 + nx}{n} = x + \dfrac{x^2}{n} \to x.
    \]
    So convergence is pointwise on $\rr$.
\end{problem}

\begin{problem}[Can limits be interchanged?]{
    \[
        \lim_{x\to a}\lim_{n\to\infty} f_n(x)
        \;\not=\;
        \lim_{n\to\infty}\lim_{x\to a} f_n(x).
    \]
}
    Define
    \[
        a_{m,n} =
        \begin{cases}
            1, & m \geq n,\\
            0, & m < n.
        \end{cases}
    \]
    Then
    \begin{align*}
        \lim_{m\to\infty}\lim_{n\to\infty} a_{m,n} &= \lim_{m\to\infty} 0 = 0,\\
        \lim_{n\to\infty}\lim_{m\to\infty} a_{m,n} &= \lim_{n\to\infty} 1 = 1.
    \end{align*}
\end{problem}

\begin{definition}[Pointwise Convergence (epsilon form)]
    $f_n \to f_0$ pointwise if for every $x\in A$ and every $\eps > 0$, there exists $N\in\n$ such that $n > N$ implies
    \[
        \abs{f_n(x) - f_0(x)} < \eps.
    \]
\end{definition}

\begin{definition}[Uniform Convergence]
    $(f_n)$ converges uniformly to $f_0$ on $A$ if for every $\eps > 0$, there exists $N\in\n$ such that for all $x\in A$, $n > N$ implies
    \[
        \abs{f_n(x) - f_0(x)} < \eps.
    \]
\end{definition}

\begin{corollary}
    Uniform convergence implies pointwise convergence, but not conversely.
\end{corollary}

\begin{theorem}{
    If $f_n \to f_0$ uniformly on $A$, and each $f_n$ is continuous, then $f_0$ is continuous.
}
    Omitted.
\end{theorem}

\begin{problem}{
    \begin{align*}
        f_n(x) &= x + \dfrac{x^2}{n},\\
        f_0(x) &= x.
    \end{align*}
}
    For $\abs{x} \leq C$,
    \[
        \abs{f_n(x) - f_0(x)} = \dfrac{x^2}{n} \leq \dfrac{C^2}{n}.
    \]
    Given $\eps > 0$, choose $N > \dfrac{C^2}{\eps}$, then for $n > N$,
    \[
        \abs{f_n(x) - f_0(x)} < \eps.
    \]
    So convergence is uniform on $[-C,C]$.
    However, on $\rr$, for $x=n$,
    \[
        \abs{f_n(n) - f_0(n)} = \dfrac{n^2}{n} = n \to \infty,
    \]
    so convergence is not uniform on $\rr$.
\end{problem}

\begin{problem}{
    $f_n(x) = x^n$ does not converge uniformly on $[0,1]$, where
    \[
        f_0(x) \coloneqq
        \begin{cases}
            0, & 0 \leq x < 1,\\
            1, & x = 1.
        \end{cases}
    \]
}
    Take $\eps = \tfrac{1}{2}$.
    For any $N\in\n$, choose $x = (1/2)^{1/N} \in (0,1)$.
    Then
    \[
        x^N = \tfrac{1}{2},
    \]
    so
    \[
        \abs{f_N(x) - f_0(x)} = \tfrac{1}{2}.
    \]
    Hence convergence is not uniform.
    However, on $[0,\gamma]$ with $\gamma < 1$, we have
    \[
        x^n \leq \gamma^n \to 0,
    \]
    so convergence is uniform on such intervals.
\end{problem}